\documentclass[10pt,a4paper]{amsart}
\usepackage[margin=19mm]{geometry}
\usepackage{amsmath,amssymb,mathtools}
\usepackage[scr=rsfs,cal=euler]{mathalfa}
\usepackage{xfrac}
\usepackage[unicode,hidelinks,bookmarks=false]{hyperref}
\newcommand{\ie}{i.e.\ }
\newcommand{\cf}{cf.\ }
\newcommand{\resp}{respectively\ }
\newcommand{\Subsection}{\subsection}
\providecommand{\nobreakdash}{\nobreak\hbox{-}}
\setlength{\emergencystretch}{2em}
\multlinegap0pt
\usepackage{bm}
\usepackage{bbm}
\usepackage{dsfont}
\usepackage{xspace}
\usepackage{cancel}
\usepackage{mathtools}
\usepackage{amsthm}
\makeatletter
\@ifundefined{theorem}{\newtheorem{theorem}{Theorem}}{}
\@ifundefined{claim}{\newtheorem{claim}{Claim}}{}
\@ifundefined{lemma}{\newtheorem{lemma}[theorem]{Lemma}}{}
\makeatother
\newcommand{\cal}{\mathcal}
\title[(Treewidth, Clique)-Boundedness and Poly-logarithmic Tree-In]{(Treewidth, Clique)-Boundedness and Poly-logarithmic Tree-Independence}
\author{Maria Chudnovsky, Ajaykrishnan E S, Daniel Lokshtanov}
\date{Source: arXiv:2510.15074v4 (CC BY 4.0)}
\begin{document}
\maketitle

\noindent\emph{Source and licence.} (Treewidth, Clique)-Boundedness and Poly-logarithmic Tree-Independence. Version arXiv:2510.15074v4 (CC BY 4.0), distributed under the Creative Commons Attribution 4.0 International licence, as recorded in the arXiv licence metadata for this exact version and shown on the abstract page https://arxiv.org/abs/2510.15074v4. Full rights notice: licenses/arxiv-2510.15074-CC-BY-4.0.txt.\par\smallskip

\noindent\textit{Editorial scope.} Counted unit: the Lemma whose source label is \texttt{lem:containers1} in the article (source0.tex lines 481--487, source label \texttt{lem:containers1}), delivered together with the article's own complete original proof of it (source0.tex lines 489--639, one \texttt{proof} environment of 151 lines). No mathematics is written, repaired or re-derived: statement and proof are the article's own text line for line. Required context retained uncounted: none. Every definition, display and label of those lines is retained unchanged. Omitted from this package and not redistributed: 1--480, 488--488, 640--1945. Nothing inside a retained excerpt is omitted or altered: every retained body is delivered exactly as the article's lines are written, so no omission receipt of any kind accompanies this package. Citations: the retained excerpts carry no citation command at all, so no bibliography entry of the article is transcribed and no reference is redistributed. Reference adaptation: none was needed and none was made; every reference inside the delivered unit resolves inside it, so no unresolved reference or citation is produced. Printed slips kept literal: no printed word, symbol, hypothesis or number of the counted statement or of its proof was altered. Layout and class: the article's own class is \texttt{amsart} with packages including geometry, graphicx, amsthm, amsmath, amssymb, todonotes, xspace, hyperref, algorithm, algpseudocode, thmtools, thm-restate, fontenc, xcolor; the wrapper is the lane's pinned \texttt{amsart} editorial wrapper at 10pt on A4, which restates the theorem-like environments the retained text needs and adds the article's own macro definitions that the retained text needs (\texttt{\textbackslash cal}), with extra wrapper packages (bm, bbm, dsfont, xspace, cancel, mathtools). The delivered numbering is the unit's own. Assets: no retained span contains a figure environment, a graphics-inclusion command, a PDF-page inclusion, a table or a TikZ picture; no image file is copied into this package, the package declares no asset dependency and no image file is redistributed with it. Licence: version arXiv:2510.15074v4 of this article is distributed under the Creative Commons Attribution 4.0 International licence, as recorded in the arXiv licence metadata for this exact version and shown on the abstract page https://arxiv.org/abs/2510.15074v4. A full rights notice is delivered at licenses/arxiv-2510.15074-CC-BY-4.0.txt. Characters: every retained excerpt is pure ASCII, so the delivered engine, XeLaTeX with its default Latin Modern text font, reports no missing character.
\begin{lemma}\label{lem:containers1}
There exists an integer $c$ with the following property.
There exists an algorithm that takes as input $(G, \omega, k, b)$ where $\omega$ and $k$ are positive integers, $b \leq \omega$ is a non-negative integer, and $G$ is a $\overline{kK_\omega}$-free graph.
The algorithm runs in time $(n+1)^{  {(2\omega \cdot \overline{\log} (n) + k+ b)^{k+b+1}}} \cdot n^c$ and outputs a
$(b, a)$-container family ${\cal F}$ such that $|{\cal F}| \leq (n+1)^{  {(2\omega \cdot \overline{\log} (n) + k+b)^{k+b+1} }}$ and $a \leq  (2 \omega \cdot \overline{\log}(n) + k+b)^{k+b+1}$.
%of vertex sets in $G$ such that $|{\cal F}| \leq (n+1)^{  {(2\omega \cdot \overline{\log} (n) + k+b)^{k+b+1} }}$, for every set $X \in {\cal F}$ we have $\alpha(X) \leq (2 \omega \cdot \overline{\log}(n) + k+b)^{k+b+1}$, and for every set $S \subseteq V(G)$ such that $\alpha(S) \leq b$ there exists an $X \in {\cal F}$ such that $S \subseteq X$.
\end{lemma} 

\begin{proof}
We begin by describing the algorithm. If $b = 0$ the algorithm outputs ${\cal F} = \{\emptyset\}$. If $b \geq 1$ and $\alpha(G) \leq \omega$ the algorithm outputs ${\cal F} = \{V(G)\}$. Suppose now that $b \geq 1$ and $\alpha(G) > \omega$.
%
Define $\rho = (1 - \frac{1}{2\omega})^{-1}$, and observe that $1 < \rho \leq 2$. The algorithm considers the set $H$ of all vertices in $G$ of degree at least $n/\rho$. 
We first check by brute force in time $n^{\omega+2}$ whether $H$ contains an independent set $I$ of size $\omega$. 

Suppose first such a set $I$ exists. Observe that $k > 1$, because otherwise $I$ is an induced $\overline{kK_\omega}$ in $G$.
%
Define $Q = \bigcap_{v \in I}N(v)$, and note that $|Q| \geq n - \sum_{v \in I} |V(G) - N(v)| \geq \frac{n}{2}$ because every vertex in $I$ has degree at least $n(1-\frac{1}{2\omega})$.
%
We have that $G[Q]$ is $\overline{(k-1)K_\omega}$-free. Indeed, suppose that $G[Q]$ contains a set $Z$ inducing a $\overline{(k-1)K_\omega}$ then $Z \cup I$ is an induced $\overline{kK_\omega}$ in $G$, a contradiction. 
%
Then, the algorithm calls itself recursively on ($G[Q]$, $\omega$, $k-1$, $b$) and obtains a family ${\cal F}_1$. 
%
It also calls itself recursively on ($G - Q$, $\omega$, $k$, $b$) and obtains a family ${\cal F}_2$.
%
The algorithm returns ${\cal F} = {\cal F}_1 \otimes {\cal F}_2$.

If no independent set $I$ of size $\omega$ exists in $H$, then $V(G) - H$ is non-empty since $\alpha(G) > \omega > \alpha(H)$.
%
The algorithm iterates through every vertex $v \notin H$ and calls itself recursively on ($G[N(v)]$, $\omega$, $k$, $b$) and obtains a family ${\cal F}^v_1$. It also calls itself recursively on  ($G - N[v]$, $\omega$, $k$, $b-1$) and obtains a family ${\cal F}^v_2$.
%
The algorithm then returns 
$${\cal F} = \{H\} \otimes\left( \bigcup_{v \notin H} \{\{v\}\} \otimes {\cal F}_1^v \otimes {\cal F}_2^v \right).$$
This completes the description of the algorithm. Each recursive call of the algorithm is made on an instance with strictly fewer vertices than $G$. Thus the algorithm always terminates and outputs a non-empty family ${\cal F}$. 
%
First we show that ${\cal F}$ covers every set with independence number at most $b$.

\begin{claim}\label{clm:family_covers1}
For every set $S \subseteq V(G)$ such that $\alpha(S) \leq b$ there exists an $X \in {\cal F}$ such that $S \subseteq X$.
\end{claim}

%\sta{\label{clm:family_covers1}
%For every set $S \subseteq V(G)$ such that $\alpha(S) \leq b$ there exists an $X \in {\cal F}$ such that $S \subseteq X$.
%}
%

\begin{proof}
We proceed by induction on $|V(G)|$. If $b=0$ then $S=\emptyset$ and the claim holds. 
Similarly, if $\alpha(G) \leq \omega$ then ${\cal F}=\{V(G)\}$ so the claim holds in this case as well.
%
We now consider the case that $b > 0$ and $\alpha(G) > \omega$.
%
Suppose first that there exists an independent set $I$ in $H$ of size $\omega$. By the induction hypothesis ${\cal F}_1$ contains a set $X_1$ that contains $S \cap Q$ and similarly ${\cal F}_2$ contains a set $X_2$ that contains $S - Q$. But then $X_1 \cup X_2 \in {\cal F}$ contains $S$. 

Suppose now that no such independent set $I$ exists. 
%If $S \subseteq H$ then every set in ${\cal F}$ contains $H$ and therefore satisfies the conditions of~(\ref{clm:family_covers1}). 
Since every set in ${\cal F}$ contains $H$ we may assume that $S - H$ is non-empty. 
% and therefore satisfies the conditions of~(\ref{clm:family_covers1}). 
Let $v$ be a vertex in $S - H$. 
%
By the induction hypothesis we have that ${\cal F}_1^v$ contains a set $X_1$ such that $S \cap N(v) \subseteq X_1$. Further, since $S - N[v]$ contains no neighbors of $v$ we have that $\alpha(S - N[v]) \leq \alpha(S) - 1 \leq b - 1$. Thus, by the induction hypothesis 
 ${\cal F}_2^v$ contains a set $X_2$ such that $S - N[v] \subseteq X_2$.
But then $S \subseteq \{v\} \cup X_1 \cup X_2$ and $\{v\} \cup X_1 \cup X_2 \in {\cal F}$, proving the claim.
%~(\ref{clm:family_covers1}). 
\end{proof}

% \medskip
% \todo[inline]{Do we need the paragraph below??}
% {\color{purple}We record two monotonicity facts used in the proofs of Claims~\ref{clm:alphaBound1}
% and~\ref{clm:FBound1}. First, for fixed $\rho > 1$ and fixed $a \geq 0$, the function
% $B_{\rho,a}$ mapping each non-negative integer $n$ to $\binom{\overline{\log}_\rho(n)+a}{a}$, is nondecreasing. Since $\binom{x+a}{a} =
% \prod_{i=1}^{a}\frac{x+i}{i}$ is a finite product of positive strictly increasing
% functions on $[0,\infty)$, it is strictly increasing there, and the claim follows since
% $\overline{\log}_\rho$ is increasing with $\overline{\log}_\rho(n) \geq 0$ for $n \geq 0$. Second, the
% function $B'_{\rho,a}$ mapping each non-negative integer $n$ to $(n+1)^{2B_{\rho,a}(n)-1}$, is nondecreasing. Indeed,
% $B_{\rho,a}(n) \geq \binom{a}{a} = 1$ for $n \geq 0$, so $2B_{\rho,a}(n)-1 \geq 1$,
% and $B_{\rho,a}$ is nondecreasing by the first fact, so for $0 \leq m \leq n$, we have
% $(m+1)^{2B_{\rho,a}(m)-1} \leq (n+1)^{2B_{\rho,a}(m)-1} \leq
% (n+1)^{2B_{\rho,a}(n)-1}$.}
Now we upper bound the independence numbers of the sets in ${\cal F}$.

\begin{claim}\label{clm:alphaBound1}
For every $X \in {\cal F}$ we have that $\alpha(X) \leq 2\omega {\overline{\log}_\rho(n) + k+b \choose k+b} - \omega$.
\end{claim}

\begin{proof}
%\sta{\label{clm:alphaBound1}
%For every $X \in {\cal F}$ we have that $\alpha(X) \leq 2\omega {\overline{\log}_\rho(n) + k+b \choose k+b} - \omega$.
%}
We proceed by induction on $n$. If $b=0$ or $\alpha(G) \leq \omega$ then we have $\alpha(X) \leq \omega \leq 2\omega {\overline{\log}_\rho(n) + k+b \choose k+b} - \omega$. Suppose now that $b \geq 1$ and $\alpha(G) > \omega$. Since $\omega \geq 1$ this implies that $n \geq \alpha(G) \geq 2$. 

If $H$ contains an independent set $I$ of size $\omega$, then by the induction hypothesis we have that for every $X_1 \in {\cal F}_1$ we have $\alpha(X_1) \leq 2\omega {\overline{\log}_\rho(n) + k-1+b \choose k-1+b} - \omega$.
%
Similarly, by the induction hypothesis (and using the fact that $|Q| \geq \frac{n}{2}$) we have that for every $X_2 \in {\cal F}_2$ we have 
$\alpha(X_2) \leq 2\omega {\overline{\log}_\rho(n)-1 + k+b \choose k+b} - \omega$.
%
It follows that for every $X \in {\cal F}$  we have 
$$\alpha(X) \leq 
2\omega {\overline{\log}_\rho(n) + k-1+b \choose k-1+b} - \omega +
2\omega {\overline{\log}_\rho(n)-1 + k+b \choose k+b} - \omega
\leq
2\omega {\overline{\log}_\rho(n) + k+b \choose k+b} - \omega.
$$

If $H$ does not contain an independent set $I$ of size $\omega$ then $\alpha(H) \leq \omega - 1$. 
%
Further, by the induction hypothesis (and using the fact that $|N(v)| < n/\rho$) we have that for every $v \notin H$ and every set $X_1 \in {\cal F}_1^v$ we have 
$\alpha(X_1) \leq 2\omega {\overline{\log}_\rho(n) - 1 + k+b \choose k+b} - \omega$. Additionally, for every $v \notin H$ and every set $X_2 \in {\cal F}_2^v$ we have $\alpha(X_2) \leq 2\omega {\overline{\log}_\rho(n) + k+b-1 \choose k+b-1} - \omega$. Thus, for every $X \in {\cal F}$  we have 
\begin{align*}
\alpha(X) & \leq \omega - 1 + 1 + 2\omega {\overline{\log}_\rho(n) - 1 + k+b \choose k+b} - \omega + 2\omega {\overline{\log}_\rho(n) + k+b-1 \choose k+b-1} - \omega \\
& \leq 2\omega {\overline{\log}_\rho(n) + k+b \choose k+b} - \omega \mbox{.}
\end{align*}
%
This concludes the proof of the claim.
%proves~(\ref{clm:alphaBound1}). 
\end{proof}

Finally, we upper bound the size of ${\cal F}$.
\begin{claim}\label{clm:FBound1}
$|{\cal F}| \leq (n+1)^{ 2 {\overline{\log}_\rho(n) + k+b \choose k+b} - 1}$.
\end{claim}

\begin{proof}
%\sta{\label{clm:FBound1}
%$|{\cal F}| \leq (n+1)^{ 2 {\overline{\log}_\rho(n) + k+b \choose k+b} - 1}$.}
The proof of this claim closely follows the proof of Claim~\ref{clm:alphaBound1}.
%
If $b=0$ or $\alpha(G) \leq \omega$ then we have $|{\cal F}| = 1 \leq n+1$. Suppose now that $b \geq 1$ and $\alpha(G) > \omega$. Since $\omega \geq 1$ this implies that $n \geq \alpha(G) \geq 2$. 

If $H$ contains an independent set $I$ of size $\omega$ then, by the induction hypothesis we have that $|{\cal F}_1| \leq (n+1)^{2 {\overline{\log}_\rho(n) + k-1+b \choose k-1+b} - 1}$.
%
Similarly, by the induction hypothesis (and using the fact that $|Q| \geq \frac{n}{2}$) we have that $|{\cal F}_2| \leq (n+1)^{2 {\overline{\log}_\rho(n)-1 + k+b \choose k+b} - 1}$.
%
It follows that
$$|{\cal F}| \leq 
(n+1)^{2 {\overline{\log}_\rho(n) + k-1+b \choose k-1+b} - 1} \cdot
(n+1)^{2 {\overline{\log}_\rho(n)-1 + k+b \choose k+b} - 1}
\leq
(n+1)^{ 2 {\overline{\log}_\rho(n) + k+b \choose k+b} - 1 }.
$$


If $H$ does not contain an independent set $I$ of size $\omega$, by the induction hypothesis (and using the fact that $|N(v)| < n/\rho$) we have that for every $v \notin H$ we have 
$|{\cal F}_1^v| \leq (n+1)^{2 {\overline{\log}_\rho(n) - 1 + k+b \choose k+b} - 1}$. Additionally, for every $v \notin H$ we have $|{\cal F}_2^v| \leq (n+1)^{2 {\overline{\log}_\rho(n) + k+b-1 \choose k+b-1} - 1}$. Hence we may conclude that  
\begin{align*}
|{\cal F}| & \leq 
n \cdot
(n+1)^{2 {\overline{\log}_\rho(n) - 1 + k+b \choose k+b} - 1} \cdot
(n+1)^{2 {\overline{\log}_\rho(n) + k+b-1 \choose k+b-1} - 1}\\
& \leq(n+1)^{ 2 {\overline{\log}_\rho(n) + k+b \choose k+b} - 1 } \mbox{.}
\end{align*}
%
This proves the claim.
%~(\ref{clm:FBound1}). 
\end{proof}

Since $\frac{1}{\log(\rho)} \leq 2\omega$, Claim~\ref{clm:FBound1} implies the claimed size bound on ${\cal F}$, while Claim~\ref{clm:alphaBound1} implies the claimed bound on the independence number of every set $X$ in ${\cal F}$. 
%
The upper bound of Claim~\ref{clm:FBound1} would apply even if duplicates of the same set in ${\cal F}$ are counted as many times as they are generated by the algorithm. Since the algorithm only spends polynomial time per set in ${\cal F}$ (counting duplicates), the running time bound follows.
\end{proof}

\end{document}
